\documentclass{article}
\usepackage{ctex}  % 加载 ctex 宏包，自动处理中文
\usepackage{amsmath, amssymb, amsfonts}
\usepackage{graphicx}
\begin{document}

\title{HMM 模型详细运行流程}
\author{}
\date{}
\maketitle

\section{概述}

隐马尔可夫模型（Hidden Markov Model, HMM）是一种用于建模时间序列数据的概率模型，主要由\textbf{状态转移矩阵}（$A$）、\textbf{观测概率矩阵}（$B$）、\textbf{初始状态概率}（$\pi$）构成。它主要用于\textbf{识别隐藏状态}，并对未来趋势进行预测。

本文件详细描述 HMM 的完整运行流程，特别强调之前讨论过的关键问题，包括 $B$ 的计算原理、前向后向概率的计算、HMM 训练的核心过程等。

\section{HMM 运行的三大阶段}

\begin{enumerate}
    \item 初始化（Initialization）
    \item 训练（Training, 使用 Baum-Welch 算法）
    \item 解码 \& 预测（Decoding \& Prediction）
\end{enumerate}

\section{初始化（Initialization）}

\subsection{隐藏状态的设定}

假设市场有两个隐藏状态：
\begin{itemize}
    \item $S_0$（上涨市场）
    \item $S_1$（下跌市场）
\end{itemize}

这些状态是不可观测的，我们只能通过股票数据（open, high, low, close, volume）来推测。

\subsection{观测数据（$O$）}

观测数据是一个 $n$ 维向量序列，即：
\[
O = \{ O_1, O_2, ..., O_T \}, \quad O_t \in \mathbb{R}^n
\]

例如，每天的市场数据可能是 5 维：
\[
O_t = [\text{open}, \text{high}, \text{low}, \text{close}, \text{volume}]
\]

\subsection{计算观测概率矩阵 $B$}

我们先随机分配一部分数据到 $S_0$（上涨），其余分配到 $S_1$（下跌），然后计算：

\begin{align}
    \mu_j &= \frac{1}{N_j} \sum_{t \in S_j} O_t \\
    \Sigma_j &= \frac{1}{N_j} \sum_{t \in S_j} (O_t - \mu_j)(O_t - \mu_j)^T
\end{align}

这些构成了\textbf{高斯分布的参数}，即：

\[
P(O_t | S_t = j) = \frac{1}{(2\pi)^{n/2} |\Sigma_j|^{1/2}} \exp \left( -\frac{1}{2} (O_t - \mu_j)^T \Sigma_j^{-1} (O_t - \mu_j) \right)
\]

最终，$B$ 由所有 $P(O_t | S_t = j)$ 组成，形成观测概率矩阵：

\[
B =
\begin{bmatrix}
P(O_1 | S_1) & P(O_2 | S_1) & \dots & P(O_T | S_1) \\
P(O_1 | S_2) & P(O_2 | S_2) & \dots & P(O_T | S_2)
\end{bmatrix}
\]

其中，$B_{ij}$ 代表在状态 $S_j$ 下，第 $i$ 个观测值的概率密度。

\subsection{假设状态转移矩阵 $A$}

状态转移矩阵描述市场状态的变化概率：

\[
A =
\begin{bmatrix}
P(S_0 \to S_0) & P(S_0 \to S_1) \\
P(S_1 \to S_0) & P(S_1 \to S_1)
\end{bmatrix}
\]

\subsection{假设初始状态分布 $\pi$}

\[
\pi = [P(S_0), P(S_1)]
\]

\section{训练（Baum-Welch 算法）}

\subsection{计算前向概率 $\alpha$（前向算法）}

\[
\alpha_t(j) = P(O_1, O_2, ..., O_t, S_t = j)
\]

\begin{align}
    \text{初始条件} (t=1): & \quad \alpha_1(j) = \pi_j \cdot B_j(O_1) \\
    \text{递归计算} (t \geq 2): & \quad \alpha_t(j) = \sum_{i} \alpha_{t-1}(i) \cdot A_{ij} \cdot B_j(O_t)
\end{align}

\subsection{计算后向概率 $\beta$（后向算法）}

\[
\beta_t(j) = P(O_{t+1}, O_{t+2}, ..., O_T | S_t = j)
\]

\begin{align}
    \text{终止条件} (t=T): & \quad \beta_T(j) = 1 \\
    \text{递归计算} (t=T-1, ..., 1): & \quad \beta_t(j) = \sum_{i} A_{ji} \cdot B_i(O_{t+1}) \cdot \beta_{t+1}(i)
\end{align}

\subsection{计算后验概率 $\gamma$（软分配）}

\[
\gamma_t(j) = \frac{\alpha_t(j) \beta_t(j)}{\sum_k \alpha_t(k) \beta_t(k)}
\]

这表示在时间 $t$，市场属于状态 $S_j$ 的概率。

\section{结论}

\begin{itemize}
    \item $B$（观测概率矩阵）通过高斯分布计算，包含均值 \& 协方差，并形成完整的矩阵。
    \item $A$（状态转移矩阵）控制状态如何变化。
    \item 前向 \& 后向算法用于计算状态概率，训练 HMM 使其拟合数据。
    \item 最终，HMM 识别市场状态，并预测趋势。
\end{itemize}

\end{document}